% Portada de M2.7: el círculo unidad con polos de un modelo AR y una señal muestreada.
\documentclass[tikz,border=0pt]{standalone}
\usepackage{resonancia}
\begin{document}
\begin{tikzpicture}
  \clip[rounded corners=6pt] (0,0) rectangle (12,3.6);
  \shade[top color=rAzul!22, bottom color=white] (0,0) rectangle (12,3.6);
  % círculo unidad con polos
  \begin{scope}[shift={(2.5,1.8)}]
    \fill[rVerde!14] (0,0) circle (1.45);
    \draw[trazo] (0,0) circle (1.45);
    \draw[eje] (-1.8,0) -- (1.9,0);
    \draw[eje] (0,-1.8) -- (0,1.9);
    \foreach \a/\r in {20/0.97, 55/0.95, 80/0.93} {
      \node[color=rNaranja] at ({1.45*\r*cos(\a)}, {1.45*\r*sin(\a)}) {$\times$};
      \node[color=rNaranja] at ({1.45*\r*cos(\a)}, {-1.45*\r*sin(\a)}) {$\times$};
    }
  \end{scope}
  % señal muestreada
  \draw[eje] (5.6,1.8) -- (11.6,1.8);
  \draw[medida, thin, smooth] plot[domain=5.7:11.4, samples=300] (\x, {1.8 + 1.1*exp(-0.14*(\x-5.7))*sin(115*(\x-5.7)*1.0)});
  \foreach \k in {0,...,17} {
    \pgfmathsetmacro{\x}{5.7+0.33*\k}
    \pgfmathsetmacro{\y}{1.8 + 1.1*exp(-0.14*(\x-5.7))*sin(115*(\x-5.7))}
    \fill[rNaranja] (\x,\y) circle (0.06);
    \draw[rNaranja, thin] (\x,1.8) -- (\x,\y);
  }
  \node[etiqueta, color=rAzul, anchor=south] at (8.6,3.0) {\small $x[k]=x(k\Delta t)$};
\end{tikzpicture}
\end{document}
